\documentclass[10pt,a4paper]{amsart}
\usepackage[margin=19mm]{geometry}
\usepackage{amsmath,amssymb,mathtools}
\usepackage[scr=rsfs,cal=euler]{mathalfa}
\usepackage{xfrac}
\usepackage[unicode,hidelinks,bookmarks=false]{hyperref}
\newcommand{\ie}{i.e.\ }
\newcommand{\cf}{cf.\ }
\newcommand{\resp}{respectively\ }
\newcommand{\Subsection}{\subsection}
\providecommand{\nobreakdash}{\nobreak\hbox{-}}
\setlength{\emergencystretch}{2em}
\multlinegap0pt
\usepackage{bm}
\usepackage{bbm}
\usepackage{dsfont}
\usepackage{xspace}
\usepackage{cancel}
\usepackage{mathtools}
\usepackage{amsthm}
\makeatletter
\@ifundefined{theorem}{\newtheorem{theorem}{Theorem}}{}
\@ifundefined{lemma}{\newtheorem{lemma}[theorem]{Lemma}}{}
\makeatother
\newcommand{\E}{\mathbb{E}}
\newcommand{\M}{{\mathcal{M}}}
\newcommand{\Pieps}{\Pi_{\M^\epsilon}}
\newcommand{\R}{\mathbb{R}}
\newcommand{\dist}{\operatorname{dist}}
\newcommand{\gap}{\Delta_\M}
\newcommand{\lse}{\operatorname{LSE}}
\newcommand{\prob}{\mathbb{P}}
\title[Diffusion Models and the Manifold Hypothesis: Log-Domain Smo]{Diffusion Models and the Manifold Hypothesis: Log-Domain Smoothing is Geometry Adaptive}
\author{Tyler Farghly, Peter Potaptchik, Samuel Howard, George Deligiannidis, Jakiw Pidstrigach}
\date{Source: arXiv:2510.02305v1 (CC BY 4.0)}
\begin{document}
\maketitle

\noindent\emph{Source and licence.} Diffusion Models and the Manifold Hypothesis: Log-Domain Smoothing is Geometry Adaptive. Version arXiv:2510.02305v1 (CC BY 4.0), distributed under the Creative Commons Attribution 4.0 International licence, as recorded in the arXiv licence metadata for this exact version and shown on the abstract page https://arxiv.org/abs/2510.02305v1. Full rights notice: licenses/arxiv-2510.02305-CC-BY-4.0.txt.\par\smallskip

\noindent\textit{Editorial scope.} Counted unit: the Lemma whose source label is \texttt{lem:tighter\_delta\_bound} in the article (source0.tex lines 1027--1035, source label \texttt{lem:tighter\_delta\_bound}), delivered together with the article's own complete original proof of it (source0.tex lines 1036--1147, one \texttt{proof} environment of 112 lines). No mathematics is written, repaired or re-derived: statement and proof are the article's own text line for line. Required context retained uncounted: lem:lse\_properties (lines 785--794); eq:lse\_properties\_1 (lines 787--789); lem:simple\_delta\_bound (lines 974--983); lem:zeta\_bound (lines 1004--1010). Every definition, display and label of those lines is retained unchanged. Omitted from this package and not redistributed: 1--784, 790--973, 984--1003, 1011--1026, 1148--1825. Nothing inside a retained excerpt is omitted or altered: every retained body is delivered exactly as the article's lines are written, so no omission receipt of any kind accompanies this package. Citations: the retained excerpts carry no citation command at all, so no bibliography entry of the article is transcribed and no reference is redistributed. Reference adaptation: none was needed and none was made; every reference inside the delivered unit resolves inside it, so no unresolved reference or citation is produced. Printed slips kept literal: no printed word, symbol, hypothesis or number of the counted statement or of its proof was altered. Layout and class: the article's own class is \texttt{article} with packages including natbib, geometry, inputenc, fontenc, hyperref, url, booktabs, amsfonts, nicefrac, microtype, xcolor, babel, latexsym, mathtools; the wrapper is the lane's pinned \texttt{amsart} editorial wrapper at 10pt on A4, which restates the theorem-like environments the retained text needs and adds the article's own macro definitions that the retained text needs (\texttt{\textbackslash E}, \texttt{\textbackslash M}, \texttt{\textbackslash Pieps}, \texttt{\textbackslash R}, \texttt{\textbackslash dist}, \texttt{\textbackslash gap}, \texttt{\textbackslash lse}, \texttt{\textbackslash prob}), with extra wrapper packages (bm, bbm, dsfont, xspace, cancel, mathtools). The delivered numbering is the unit's own. Assets: no retained span contains a figure environment, a graphics-inclusion command, a PDF-page inclusion, a table or a TikZ picture; no image file is copied into this package, the package declares no asset dependency and no image file is redistributed with it. Licence: version arXiv:2510.02305v1 of this article is distributed under the Creative Commons Attribution 4.0 International licence, as recorded in the arXiv licence metadata for this exact version and shown on the abstract page https://arxiv.org/abs/2510.02305v1. A full rights notice is delivered at licenses/arxiv-2510.02305-CC-BY-4.0.txt. Characters: every retained excerpt is pure ASCII, so the delivered engine, XeLaTeX with its default Latin Modern text font, reports no missing character.
\begin{lemma}\label{lem:lse_properties}
For any \(\{x_i\}_{i=1}^N \subset \R\) and \(\{\varepsilon_i\}_{i=1}^N \subset \R\), we have
\begin{equation}\label{eq:lse_properties_1}
    \lse(\{x_i + \epsilon_i\}_{i=1}^N) - \lse(\{x_i\}_{i=1}^N) = \int_0^1 \frac{\sum_{i=1}^N \exp(x_i + r \epsilon_i) \epsilon_i}{\sum_{i=1}^N \exp(x_i + r \epsilon_i)} dr.
\end{equation}
In particular, we have that
\begin{equation}\label{eq:lse_properties_2}
    \lse(\{x_i + \epsilon_i\}_{i=1}^N) - \lse(\{x_i\}_{i=1}^N) \leq \max\{\epsilon_i\}.
\end{equation}
\end{lemma}

\begin{lemma}\label{lem:simple_delta_bound}
Suppose that \(\gap := \dist(\{x_i\}_{i=1}^N, \M) < \infty\) and \(\tau_\M > 0\). Then, for any \(x \in \R^d, i \in [N]\) we have,
\begin{equation*}
    |\Delta_i| \leq \mu_\epsilon |\zeta| \Big ( \Big ( \tfrac{1}{2 \tau_\M} d_i^2 \Big ) \wedge d_i + \gap \Big ), \qquad d_i = \| \Pi_\M(Y/\mu_\epsilon) - \Pi_\M(x_i)\|,
\end{equation*}
where we define the quantity,
\begin{equation*}
    \zeta := \|Y - \Pieps(Y)\| - \E \Big [ \|Y - \Pieps(Y)\|^2 \Big ]^{1/2}.
\end{equation*}
\end{lemma}

\begin{lemma}\label{lem:zeta_bound}
Let \(\zeta\) be as in Lemma \ref{lem:simple_delta_bound} and suppose that \(K_\M, K_{\max, \M} < \infty\), then we have that,
\begin{equation*}
    \E[|\zeta|^2|S]^{1/2} \leq 2 K_\M, \qquad |\zeta| \leq K^{\max}_\M + K_\M,
\end{equation*}
almost surely.
\end{lemma}

\begin{lemma}\label{lem:tighter_delta_bound}
Consider the setting of Lemma \ref{lem:simple_delta_bound}, then for any \(x \in \R^d\), it holds that
\begin{align*}
    &\E[|\Delta \lse_\M(x)|S]\\
    & \qquad \leq \frac{8 K_\M}{\tau_\M} \bigg ( 1 + 2\log \big ( \tfrac{\tau_\M}{K_\M} \big )_+ + (\gap + \tau_\M + K_\M + D_x ) \frac{12\gap}{\sigma_\epsilon^2} + \Big ( 5 K_{\max, \M}^{1/2} K_\M^{1/2} + D_x \Big )^2 \frac{\prob(\mathcal{E}_x|S)^{1/2}}{\sigma_\epsilon^2}\\
    & \qquad \qquad + 2\inf_{\varepsilon_0 > 0} \bigg \{ \log \Big ( \E_{Y \sim k_x} [ \hat{\mu}_{\operatorname{data}}(B_{\varepsilon_0}(\Pi_\M(Y/\mu_\epsilon)))^{-1} | S ] \Big )_+ + \Big ( 1 + \tfrac{K_\M + D_x}{\tau_\M} \Big ) \frac{\varepsilon_0^2}{\sigma_\epsilon^2} \bigg \} \bigg ),
\end{align*}
for some universal constant \(\tilde{C} > 0\), where we define \(\mathcal{E}_x = \{ K_\M + D_x + 2 |\zeta| \geq \mu_\epsilon \tau_\M / 2 \}\), \(D_x = \|x - \Pieps(x)\|\).
\end{lemma}

\begin{proof}
For this bound, we begin with \eqref{eq:lse_properties_1} of Lemma \ref{lem:lse_properties} to obtain
\begin{equation*}
    | \Delta \lse_\M(x) | \leq \int_0^1 \frac{\sum_{i \in [N]} \exp(-\|\tilde{Y} - \mu_\epsilon x_i\|^2/2\sigma_\epsilon^2-r\Delta_i/2\sigma_\epsilon^2) |\Delta_i|/2\sigma_\epsilon^2}{\sum_{i \in [N]} \exp(-\|\tilde{Y} - \mu_\epsilon x_i\|^2/2\sigma_\epsilon^2-r\Delta_i/2\sigma_\epsilon^2)} dr.
\end{equation*}
We decompose this further as,
\begin{align*}
    | \Delta \lse_\M(x) | &\leq \int_0^1 \frac{\sum_{i \in I_1} \exp(-\|\tilde{Y} - \mu_\epsilon x_i\|^2/2\sigma_\epsilon^2-r\Delta_i/2\sigma_\epsilon^2) |\Delta_i|/2\sigma_\epsilon^2}{\sum_{i \in I_1} \exp(-\|\tilde{Y} - \mu_\epsilon x_i\|^2/2\sigma_\epsilon^2-r\Delta_i/2\sigma_\epsilon^2)} dr\\
    & \qquad + \int_0^1 \frac{\sum_{i \in I_1^\complement} \exp(-\|\tilde{Y} - \mu_\epsilon x_i\|^22\sigma_\epsilon^2-r\Delta_i/2\sigma_\epsilon^2) |\Delta_i|/2\sigma_\epsilon^2}{\sum_{i \in I_0} \exp(-\|\tilde{Y} - \mu_\epsilon x_i\|^2/2\sigma_\epsilon^2-r\Delta_i/2\sigma_\epsilon^2)} dr \\
    &\leq \mathbbm{1}_{|I_1| > 0} \max_{i \in I_{1}} \bigg \{ \frac{|\Delta_i|}{2\sigma_\epsilon^2} \bigg \} + \int_0^1 \frac{\sum_{i \in I_{1}^\complement} \exp(-\|\tilde{Y} - \mu_\epsilon x_i\|^2/2\sigma_\epsilon^2-r\Delta_i/2\sigma_\epsilon^2) |\Delta_i|/2\sigma_\epsilon^2}{\sum_{i \in I_0} \exp(-\|\tilde{Y} - \mu_\epsilon x_i\|^2/2\sigma_\epsilon^2-r\Delta_i/2\sigma_\epsilon^2)} dr \\
    &=: \mathtt{A} + \mathtt{B},
\end{align*}
where we define the sets,
\begin{gather*}
    I_0 = \{i \in [N]: \|\Pi_\M(Y/\mu_\epsilon) - \Pi_\M(x_i)\| \leq \varepsilon_0\}, \qquad I_1 = \{i \in [N]: \|\Pi_\M(Y/\mu_\epsilon) - \Pi_\M(x_i)\| \leq \varepsilon_1\}
\end{gather*}
for some random quantities \(\varepsilon_0, \varepsilon_1 > 0\).

The quantity \(\mathtt{A}\) can be bounded directly using Lemma \ref{lem:simple_delta_bound}. From this, we obtain,
\begin{align*}
    \mathtt{A} \leq \frac{\mu_\epsilon}{\sigma^2_\epsilon} |\zeta| \Big ( \tfrac{1}{2 \tau_\M} \varepsilon_1^2 + \gap \Big ).
\end{align*}
To bound \(\mathtt{B}\), we proceed with the following upper bound,
\begin{align}
    \mathtt{B} &\leq \frac{\sum_{j \in I_1^\complement} \exp(-\|\tilde{Y} - \mu_\epsilon x_j\|/2\sigma_\epsilon^2 + |\Delta_j|/2\sigma_\epsilon^2) |\Delta_j|/2\sigma_\epsilon^2}{\sum_{i \in I_0}\exp(-\|\tilde{Y} - \mu_\epsilon x_i\|/2\sigma_\epsilon^2-|\Delta_i|/2\sigma_\epsilon^2)} \nonumber\\
    &\leq |I_0|^{-1} \max_{i \in I_0}\sum_{j \in I_1^\complement} \exp(\|\tilde{Y} - \mu_\epsilon x_i\|^2/2\sigma_\epsilon^2 - \|\tilde{Y} - \mu_\epsilon x_j\|^2/2\sigma_\epsilon^2 + |\Delta_i|/2\sigma_\epsilon^2 + |\Delta_j|/\sigma_\epsilon^2),\label{eq:drbg_3}
\end{align}
where we have used the fact that \(r \leq \exp(r)\). To further control \(\mathtt{B}\), we control the quantity inside of the exponential function by choosing \(\varepsilon_0\) sufficiently small and \(\varepsilon_1\) sufficiently large so that the quantity in the exponential becomes negative. This allows for control of \(\mathtt{B}\) by taking \(\epsilon\) sufficiently small.

We start by controlling \(\|\tilde{Y} - \mu_\epsilon x_i\|^2-\|\tilde{Y} - \mu_\epsilon x_j\|^2\). For any \(i \in I_0, j \in [N]\), we have that
\begin{align}
    &\|\tilde{Y} - \mu_\epsilon x_i\|^2-\|\tilde{Y} - \mu_\epsilon x_j\|^2\nonumber\\
    & \qquad = \| \Pieps(Y) - \mu_\epsilon x_i\|^2  + 2 \langle \tilde{Y} - \Pieps(Y), \Pieps(Y) - \mu_\epsilon x_i \rangle - \|\Pieps(Y) - \mu_\epsilon x_j\|^2\nonumber\\
    & \qquad \qquad - 2 \langle \tilde{Y} - \Pieps(Y), \Pieps(Y) - \mu_\epsilon x_j \rangle.\label{eq:drbg_3h}
\end{align}
The first term is bounded using the fact that \(\|\Pieps(Y) - \mu_\epsilon x_i\| \leq \mu_\epsilon \varepsilon_0 + \mu_\epsilon \gap\). The second term is controlled using the technique from the proof of Lemma \ref{lem:simple_delta_bound} to deduce the bound,
\begin{align}
    &\langle \tilde{Y} - \Pieps(Y), \Pieps(Y) - \mu_\epsilon x_i \rangle\\
    & \qquad \leq \mu_\epsilon \|\tilde{Y} - \Pieps(Y)\| \Big ( \big ( \tfrac{1}{2\tau_\M} \|\Pi_\M(Y/\mu_\epsilon) - \Pi_\M(x_i)\|^2 \big ) \wedge \|\Pi_\M(Y/\mu_\epsilon) - \Pi_\M(x_i)\| + \gap \Big )\nonumber\\
    & \qquad \leq \mu_\epsilon (K_\M + D_x) \Big ( \big ( \tfrac{1}{2\tau_\M} \varepsilon_0^2 \big ) \wedge \varepsilon_0 + \gap \Big ),\label{eq:drbg_3hh}
\end{align}
where in the second line, we use that
\begin{align*}
    \|\tilde{Y} - \Pieps(Y)\| &= \E[\|Y - \Pieps(Y)\|^2]^{1/2} \\
    &\leq \E \Big [ \big | \|Y - \Pieps(Y)\| - \|x - \Pieps(x)\| \big |^2 \Big ]^{1/2} + \|x - \Pieps(x)\|\\
    &\leq K_\M + D_x.
\end{align*}
Similarly, the fourth term of \eqref{eq:drbg_3h} is bounded by,
\begin{align}
    &-\langle \tilde{Y} - \Pieps(Y), \Pieps(Y) - \mu_\epsilon x_j \rangle \nonumber\\
    & \qquad \leq \mu_\epsilon (K_\M + D_x) \Big ( \big ( \tfrac{1}{2\tau_\M} \|\Pi_\M(Y/\mu_\epsilon) - \Pi_\M(x_j)\|^2 \big ) \wedge \|\Pi_\M(Y/\mu_\epsilon) - \Pi_\M(x_j)\| + \gap \Big ),\label{eq:drbg_4h}
\end{align}
and finally, the third term of \eqref{eq:drbg_3h} is controlled using Young's inequality to obtain,
\begin{align}
    \|\Pieps(Y) - \mu_\epsilon x_j\|^2 &= \|\Pieps(Y) - \Pieps(\mu_\epsilon x_j)\|^2 + \|\Pieps(\mu_\epsilon x_j) - \mu_\epsilon x_j\|^2 \nonumber\\
    & \qquad + 2 \langle \Pieps(Y) - \Pieps(\mu_\epsilon x_j), \Pieps(\mu_\epsilon x_j) - \mu_\epsilon x_j \rangle \nonumber\\
    &\geq \frac{3}{4}\|\Pieps(Y) - \Pieps(\mu_\epsilon x_j)\|^2 - 3 \|\Pieps(\mu_\epsilon x_j) - \mu_\epsilon x_j\|^2 \nonumber\\
    &\geq \frac{3}{4}\|\Pieps(Y) - \Pieps(\mu_\epsilon x_j)\|^2 - 3 \mu_\epsilon^2 \gap^2,\label{eq:drbg_4}
\end{align}
Thus, substituting \eqref{eq:drbg_3hh}, \eqref{eq:drbg_4h} and \eqref{eq:drbg_4} in to \eqref{eq:drbg_3h} leads to the bound,
\begin{align}
    &\|\tilde{Y} - \mu_\epsilon x_i\|^2-\|\tilde{Y} - \mu_\epsilon x_j\|^2 \nonumber\\
    & \qquad \leq \mu_\epsilon^2 (\varepsilon_0 + \gap)^2 + 2\mu_\epsilon (K_\M + D_x) \Big ( \big ( \tfrac{1}{2\tau_\M} \varepsilon_0^2 \big ) \wedge \varepsilon_0 + \gap \Big ) - \frac{3}{4} \|\Pieps(Y) - \mu_\epsilon \Pi_\M(x_j)\|^2 + 3 \mu_\epsilon^2 \gap^2 \nonumber\\
    & \qquad \qquad + 2 \mu_\epsilon (K_\M + D_x) \Big ( \big ( \tfrac{1}{2\tau_\M}\|\Pi_\M(Y/\mu_\epsilon) - \mu_\epsilon \Pi_\M(x_j)\|^2 \big ) \wedge \|\Pi_\M(Y/\mu_\epsilon) - \Pi_\M(x_j)\| + \gap \Big ) \nonumber\\
    & \qquad \leq \mu_\epsilon \Big ( 2 \mu_\epsilon + \tfrac{1}{\tau_\M} K_\M + \tfrac{1}{\tau_\M} D_x \Big ) \varepsilon_0^2 + \mu_\epsilon \Big ( 5 \mu_\epsilon \gap + 4 K_\M + 4 D_x \Big ) \gap - \tfrac{3}{4} \|\Pieps(Y) - \mu_\epsilon \Pi_\M(x_j)\|^2 \nonumber \\
    & \qquad \qquad + 2 ( K_\M + D_x) \big ( \tfrac{1}{2 \mu_\epsilon \tau_\M}\|\Pieps(Y) - \mu_\epsilon \Pi_\M(x_j)\|^2 \big ) \wedge \|\Pieps(Y) - \mu_\epsilon \Pi_\M(x_j)\|. \label{eq:drbg_44}
\end{align}
Continuing with bounding the contents of the exponential function in \eqref{eq:drbg_3}, we next control \(|\Delta_i| + 2 |\Delta_j|\), using Lemma \ref{lem:simple_delta_bound} to obtain,
\begin{equation}
    |\Delta_i| + 2 |\Delta_j| \leq 2 \mu_\epsilon |\zeta| \Big ( \tfrac{1}{2\tau_\M} \varepsilon_0^2 + 2 \big ( \tfrac{1}{2\tau_\M} \|\Pi_\M(Y/\mu_\epsilon) - \Pi_\M(x_j) \|^2 \big ) \wedge \|\Pi_\M(Y/\mu_\epsilon) - \Pi_\M(x_j) \|  + 3 \gap \Big ).\label{eq:drbg_5}
\end{equation}
Therefore, combining \eqref{eq:drbg_44} and \eqref{eq:drbg_5}, we obtain the bound,
\begin{align*}
    &\|\tilde{Y} - \mu_\epsilon x_i\|^2-\|\tilde{Y} - \mu_\epsilon x_j\|^2 + |\Delta_i| + 2|\Delta_j|\\
    & \leq \mu_\epsilon \Big ( 2\mu_\epsilon + \tfrac{1}{\tau_\M} (|\zeta| + K_\M + D_x) \Big ) \varepsilon_0^2 + \mu_\epsilon \Big ( 5 \mu_\epsilon \gap + 6 |\zeta| + 4 K_\M + 4 D_x \Big ) \gap - \tfrac{3}{4} \|\Pieps(Y) - \mu_\epsilon \Pi_\M(x_j)\|^2\\
    & \qquad + 2 (K_\M + D_x + 2 |\zeta|) \big ( \tfrac{1}{2 \mu_\epsilon \tau_\M}\|\Pieps(Y) - \mu_\epsilon \Pi_\M(x_j)\|^2 \big ) \wedge \|\Pieps(Y) - \mu_\epsilon \Pi_\M(x_j)\|\\
    & \leq \mu_\epsilon \Big ( 2\mu_\epsilon + \tfrac{1}{\tau_\M} (|\zeta| + K_\M + D_x) \Big ) \varepsilon_0^2 + \mu_\epsilon \Big ( 5 \mu_\epsilon \gap + 6 |\zeta| + 4 K_\M + 4 D_x \Big ) \gap - \tfrac{1}{2} \|\Pieps(Y) - \mu_\epsilon \Pi_\M(x_j)\|^2\\
    & \qquad + 2 \mathbbm{1}_{\mathcal{E}_x} (K_\M + D_x + 2 |\zeta|) \|\Pieps(Y) - \mu_\epsilon \Pi_\M(x_j)\|,
\end{align*}
where the indicator function contains the event \(\mathcal{E}_x = \{ K_\M + D_x + 2 |\zeta| \geq \mu_\epsilon \tau_\M / 4 \}\).

Using this bound, we choose a value of \(\varepsilon_1\) that guarantees that the contents of the exponential function in \eqref{eq:drbg_3} is negative.
By solving the quadratic, it follows that to have \(\|\tilde{Y} - \mu_\epsilon x_i\|^2-\|\tilde{Y} - \mu_\epsilon x_j\|^2 + |\Delta_i| + 2|\Delta_j| \leq - \mu_\epsilon \kappa\), for some \(\kappa > 0\), it is sufficient to have \(\|\Pieps(Y) - \Pieps(\mu_\epsilon x_j)\| \geq \mu_\epsilon \varepsilon_1\) with,
\begin{align}
    \mu_\epsilon^2 \varepsilon_1^2 &= 6 \mu_\epsilon \kappa + 6 \mu_\epsilon \Big ( 2\mu_\epsilon + \tfrac{1}{\tau_\M} (|\zeta| + K_\M + D_x) \Big ) \varepsilon_0^2 + 6 \mu_\epsilon \Big ( 5 \mu_\epsilon \gap + 6 |\zeta| + 4 K_\M + 4 D_x \Big ) \gap \nonumber\\
    & \qquad + 8 \mathbbm{1}_{\tilde{\mathcal{E}}_x} (K_\M + D_x + 2 |\zeta|)^2.\label{eq:drbg_2}
\end{align}

Substituting this into \eqref{eq:drbg_3}, we then obtain the following bound for \(\mathtt{B}\),
\begin{equation*}
    \E[\mathtt{B}|S] \leq \E \bigg [ \frac{|I_1^\complement|}{|I_0|} \bigg | S \bigg ] \exp(-\mu_\epsilon\kappa/2\sigma^2_\epsilon).
\end{equation*}
We now return to bounding \(\mathtt{A}\) with this choice of \(\varepsilon_1\)
to obtain that,
\begin{align*}
    \E[\mathtt{A}|S] &\leq \frac{\mu_\epsilon K_\M}{\sigma^2_\epsilon} \bigg ( \frac{1}{\tau_\M \mu_\epsilon^2} \bigg ( 6 \mu_\epsilon \kappa + 6 \mu_\epsilon \Big ( 2\mu_\epsilon + \tfrac{3}{\tau_\M} K_\M + \tfrac{1}{\tau_\M} D_x \Big ) \varepsilon_0^2 + 6 \mu_\epsilon \Big ( 5 \mu_\epsilon \gap + 16 K_\M + 4 D_x \Big ) \gap \\
    & \qquad + 8 \prob(\mathcal{E}_x|S)^{1/2} (5 K_{\max, \M}^{1/2} K_\M^{1/2} + D_x)^2 \bigg ) + 2 \gap \bigg ),
\end{align*}
where we utilise the bounds in Lemma \ref{lem:zeta_bound} to control \(|\zeta|\).
We then optimise \(\kappa\) by choosing,
\begin{equation*}
    \kappa = \frac{2 \sigma_\epsilon^2}{\mu_\epsilon} \log \bigg ( \E \Big [ |I_1^\complement|/|I_0| \Big | S \Big ] \frac{\tau_\M}{K_\M} \bigg )_+,
\end{equation*}
which produces the bound,
\begin{align*}
    \E[|\Delta \lse(x)|S] &\leq \frac{K_\M}{\tau_\M} \bigg ( 1 + 12 \log \Big ( \E \Big [ |I_1^\complement|/|I_0| \Big | S \Big ] \tfrac{\tau_\M}{K_\M} \Big )_+ + \Big ( 2 + \tfrac{3}{\tau_\M} K_\M + \tfrac{1}{\tau_\M} D_x \Big ) \frac{6\varepsilon_0^2}{\sigma_\epsilon^2} \\
    & \quad + \Big ( 5 \gap + 16 K_\M + 4 D_x \Big ) \frac{6 \gap}{\sigma_\epsilon^2} +\frac{8 \prob(\mathcal{E}_x|S)^{1/2}}{\sigma_\epsilon^2} (5 K_{\max, \M}^{1/2} K_\M^{1/2} + D_x)^2 \bigg ) \bigg ) + \frac{2 K_\M \gap}{\sigma_\epsilon^2},
\end{align*}
where we have used that \(\mu_\epsilon \leq 1\) to simplify the expression.

To conclude the proof, we use the fact that \(|I_0| = N \hat{\mu}_{\operatorname{data}}(B_{\varepsilon_0}(\Pi_\M(Y/\mu_\epsilon))\) and \(|I_1^\complement| \leq N\). Then, optimising over \(\varepsilon_0\) leads to the bound in the statement.
\end{proof}

\end{document}
